\section{Related Work}

\paragraph{Data scale for robot policies}
RT-2~\cite{rt2} co-trains a vision-language backbone on web data and robot trajectories; OpenVLA~\cite{openvla} and $\pi_0$~\cite{pi0} follow the same recipe on the Open X-Embodiment collection~\cite{oxe}, which pools demonstrations from many labs to reach a scale no single lab could record. Egocentric human video is the other large source. Ego4D~\cite{ego4d} and EPIC-KITCHENS~\cite{epickitchens100} show what thousands of hours of first-person household footage look like, and several VLA efforts now pretrain on human video before fine-tuning on robot data. Our setting is the opposite end of that scale: a few recordings from one environment, with the question of how far they can be stretched before a robot is ever trained.

\paragraph{Augmentation for robot learning}
Domain randomization~\cite{tobin2017} randomizes textures, lighting and camera pose in simulation so that the real world looks like one more sample; Peng et al.~\cite{peng2018} extend the idea to dynamics. The approach is deliberately agnostic to semantics, and that is its limit for appearance augmentation of real video: there is no mechanism to keep a variation within the set a household actually produces. ROSIE~\cite{rosie} and GenAug~\cite{genaug} close that gap for images by prompting a text-to-image diffusion model with a language description of a plausible change and inpainting it into robot observations. We borrow the premise that a language model should decide what to vary, and we move it to video, where the edit has to hold across 60 frames of a moving hand, and we make the decision procedure explicit as a schema rather than a free-form prompt.

\paragraph{Structure-conditioned video generation}
ControlNet~\cite{controlnet} showed that a frozen image diffusion model can be steered by edge, depth or segmentation maps through a trainable side branch. Gen-1~\cite{gen1} brought the same control to video by conditioning on per-frame depth while a text prompt supplies appearance. NVIDIA's Cosmos platform~\cite{cosmos2025} trains video diffusion models for physical AI and releases three families: Predict for forecasting future frames~\cite{cosmospredict2}, Transfer for multi-control video-to-video adaptation~\cite{cosmostransfer1}, and Reason for scene understanding~\cite{cosmosreason1}. Cosmos Transfer adds adaptive per-control weights so that edge, depth, segmentation and blurred-RGB branches can be mixed. We use Transfer~2.5~\cite{cosmos25,cosmostransfer25} as the renderer and report the weights that work for egocentric manipulation, which the model card does not cover. We also use the newer Cosmos~3 Nano~\cite{cosmos3,cosmos3nano} as the reasoner, since its reasoning tower returns bounding boxes with its text, and we found its generation tower did not yet support video-to-video transfer.

\paragraph{Segmentation and preprocessing}
Our masks come from SAM2~\cite{sam2}, which propagates a box prompt through a video and returns temporally consistent masks; the subtraction step in Sec.~\ref{sec:mask} is a set operation on those masks. Clip boundaries are detected with TransNetV2~\cite{transnetv2}, and the edge and depth controls use the Canny detector~\cite{canny1986} and a monocular depth network~\cite{depthanything}, both as packaged with the Cosmos Transfer release.
